\documentclass[11pt]{amsart}
\usepackage{amsmath,amssymb,amsthm,booktabs}
\usepackage[margin=1.15in]{geometry}
\usepackage{hyperref}
\newtheorem{theorem}{Theorem}
\newtheorem{proposition}[theorem]{Proposition}
\newtheorem{lemma}[theorem]{Lemma}
\newtheorem{corollary}[theorem]{Corollary}
\theoremstyle{definition}
\newtheorem{example}[theorem]{Example}
\newtheorem{remark}[theorem]{Remark}
\newcommand{\Prob}{\mathbb{P}}
\newcommand{\E}{\mathbb{E}}
\newcommand{\eps}{\varepsilon}
\newcommand{\cF}{\mathcal{F}}
\newcommand{\cK}{\mathcal{K}}
\newcommand{\cR}{\mathcal{R}}
\newcommand{\cT}{\mathcal{T}}

\title[Sublinear preferential attachment with fixed out-degree]{The local limit of sublinear preferential
attachment graphs with fixed out-degree}
\author{Yitzchak Shmalo}
\date{November 2026}

\begin{document}
\begin{abstract}
We consider the preferential attachment graph in which every new vertex sends $m\ge2$ edges to the existing
vertices, the endpoints being chosen one after another with probability proportional to $f(\text{in-degree})$, where
the attachment function $f$ is bounded away from zero, has bounded increments and satisfies $f(k)/k\to0$. No
monotonicity or concavity is assumed. We prove that the total weight of the first $n$ vertices divided by $n$
converges almost surely to the unique solution $c_f$ of an explicit scalar equation, that the empirical in-degree
distribution converges almost surely to the law of a pure birth process with rates $mf(k)/c_f$ observed at an
independent exponential time, and that the graph converges locally in probability to an infinite random tree. The
probability of every finite ball in the limit is given by an explicit integral over the ages of its vertices. This
answers, for this class of attachment functions, a question of Deijfen and van der Hofstad on local limits beyond
the affine case when the out-degree is fixed.
\end{abstract}
\maketitle

\section{Introduction}

\subsection{The model and the question}

Fix an integer $m\ge2$ and a function $f:\{0,1,2,\dots\}\to(0,\infty)$. The graph $G_1$ consists of the vertex $1$
and no edge. Given $G_t$, whose vertex set is $[t]=\{1,\dots,t\}$, the graph $G_{t+1}$ is obtained by adding the
vertex $t+1$ and $m$ directed edges from $t+1$ to vertices of $[t]$. The endpoints are chosen one after another.
If $D_i$ denotes the current in-degree of $i\in[t]$, which includes the edges already sent by $t+1$, the next
endpoint is $i$ with probability
\[
     \frac{f(D_i)}{\sum_{j\le t}f(D_j)} .
\]
Thus the in-degrees are updated after every edge, multiple edges may occur, and there are no loops. We write
$D_i(n)$ for the in-degree of $i$ in $G_n$ and $S_n=\sum_{i\le n}f(D_i(n))$ for the total weight.

When $f(k)=k+\delta$ is affine, the weights form a P\'olya urn, the graph has an exchangeable representation,
and Berger, Borgs, Chayes and Saberi~\cite{BBCS} used it to identify the local limit, the P\'olya point tree;
see \cite{GHHR} and \cite[Chapter~5]{vdH2} for extensions, including random out-degrees. For $m=1$ the graph is a
tree, which embeds in a continuous-time branching process for every $f$, and Rudas, T\'oth and
Valk\'o~\cite{RTV} obtained its local limit in this way. Dereich and M\"orters~\cite{DM1,DM2} treat concave
attachment functions in a model in which a new vertex connects to each old vertex independently, so
that the out-degree is random. Open Problem~24 of~\cite{Aletti}, contributed by Deijfen and van der Hofstad,
asks for the local limit when the out-degree is fixed, $m\ge2$, and $f$ is not affine: the branching process
embedding does not extend to this case, and the P\'olya urn representation is not available. Banerjee, Deka and
Olvera-Cravioto~\cite{BDO} obtain local limits for graphs constructed by collapsing the vertices of a tree with a
general attachment function; in a collapsed graph the weight of a vertex is the sum of the weights of the tree
vertices it is made of, which is a function of its degree only when $f$ is affine.

\subsection{The result}

We assume
\begin{itemize}
\item[\textup{(E1)}] $f(k)\ge\alpha$ for all $k$, for some $\alpha>0$;
\item[\textup{(E2)}] $|f(k+1)-f(k)|\le L$ for all $k$, for some $L<\infty$;
\item[\textup{(E3)}] $f(k)/k\to0$ as $k\to\infty$.
\end{itemize}
Examples are all functions with values in a compact subset of $(0,\infty)$, the functions $(k+\delta)^a$ with
$\delta>0$ and $0<a<1$, and $1+\log(k+1)$. For $d>0$ let $p(d)=(p_k(d))_{k\ge0}$ be the probability vector
\begin{equation}\label{eq:pk}
   p_k(d)=\frac{d}{d+mf(k)}\prod_{j=0}^{k-1}\frac{mf(j)}{d+mf(j)} ,
\end{equation}
which is the law of $X(T)$ when $X$ is the pure birth process started at $0$ with rates $(m/d)f(k)$ and $T$ is
an independent exponential time of mean one.

For a rooted graph $(G,u)$ and $r\ge0$, let $B_r(G,u)$ be the subgraph induced by the vertices at distance at most
$r$ from $u$, where the distance ignores the orientation of the edges; $B_r(G,u)$ is rooted at $u$ and keeps the
orientations. We write $\cong$ for isomorphism of rooted directed multigraphs. A sequence of finite random graphs
$G_n$ \emph{converges locally in probability} to a random rooted graph $(\cT,o)$ if, for every $r$ and every finite
rooted directed multigraph $H$,
\[
    \frac1{|G_n|}\,\#\bigl\{u\in G_n:\ B_r(G_n,u)\cong H\bigr\}\ \longrightarrow\ \Prob\bigl(B_r(\cT,o)\cong H\bigr)
\]
in probability \cite{BS,vdH2}.

\begin{theorem}\label{thm:main}
Assume \textup{(E1)--(E3)}.
\begin{enumerate}
\item[\textup{(i)}] There is a unique $c_f>0$ such that
\[
     \sum_{k\ge1}\ \prod_{j=0}^{k-1}\frac{mf(j)}{c_f+mf(j)}=m ,
\]
that is, such that $p(c_f)$ has mean $m$.
\item[\textup{(ii)}] Almost surely $S_n/n\to c_f$, and the proportion of vertices of $G_n$ with in-degree $k$
converges to $p_k(c_f)$ for every $k$, with convergence of the sequences in $\ell^1$. Moreover
$\max_{i\le n}D_i(n)/n\to0$ in probability.
\item[\textup{(iii)}] $G_n$ converges locally in probability to a random rooted directed tree $(\cT_f,o)$. This
tree is infinite and has no multiple edges, every vertex has out-degree $m$, the in-degree of the root has law
$p(c_f)$, and for every admissible ball shape $H$ of radius $r$ the probability $\Prob(B_r(\cT_f,o)\cong H)$ is the
integral $\mathbf P_H$ given in \eqref{eq:PH} and Proposition~\ref{prop:factors}, with $c=c_f$.
\end{enumerate}
\end{theorem}

Constant $f$ is included and gives uniform attachment, for which $c_f=f(0)$ and $p(c_f)$ is geometric. The
theorem does not cover attachment functions with $f(k)/k\to\gamma\in(0,\infty)$, affine or not, functions with
unbounded increments, or functions with $\inf f=0$. Other conventions for the $m$
edges of a vertex, such as in-degrees that are updated only after the $m$ edges, are not treated either.

\begin{remark}[The limiting tree]\label{rem:tree}
The formula for $\mathbf P_H$ has the following reading. Let $\lambda=m/c_f$ and $h_\lambda(u)=\E f(X(u))$, where
$X$ is the pure birth process started at $0$ with rates $\lambda f(k)$. Every vertex of $\cT_f$ has an age in
$(0,1)$, and the age of the root is uniform. A vertex of age $x$ sends its $m$ edges to parents whose ages are
independent with density $h_\lambda(\log(x/y))/(c_fx)$ on $(0,x)$; for a vertex reached from one of its parents
this applies to the other $m-1$ edges. The density is uniform only when $f$ is constant. The children of the root, and of any vertex reached from one of its parents, are born at the jump
times of a birth process which starts from $0$ at the age $x$ of the vertex and jumps from $k$ to $k+1$ at rate
$\lambda f(k)/t$ at time $t\in(x,1]$. For a vertex of age $y$ reached from one of its children, of age $x$, the
law of this birth process is weighted by $f$ evaluated at its state just before time $x$, and a jump is added at
time $x$. For $f\equiv1$ this is the local limit of uniform attachment.
\end{remark}

\subsection{The proof}

Without an urn representation the first difficulty is the normalisation $S_n$. We write the empirical in-degree
distribution as a stochastic approximation in $\ell^1$ and identify its limit through a time change that turns
the limiting differential equation into a linear one; the linear equation is solved by a renewal equation, and the
key renewal theorem gives convergence to $p(c_f)$ uniformly over the compact set of probability vectors with mean
at most $m$ (Section~\ref{sec:normal}). No monotonicity of $f$ is needed, and the equilibrium need not be stable
in any linearised sense that we would have to verify. The local limit is then obtained from an exact change of
measure: the probability that prescribed slots choose prescribed endpoints equals an expectation, for the
process in which these choices are forced, of a product of ratios $f(D)/S$ (Section~\ref{sec:local}). Finitely
many forced choices do not change the normalisation, the in-degrees of finitely many tagged vertices become
independent birth processes in the logarithm of the age, and a count of embeddings with a second moment
computation for two independent roots completes the proof. The a priori bounds on the size of balls and on old
vertices are in Section~\ref{sec:apriori}. Sections~\ref{sec:apriori}--\ref{sec:local} are written for bounded
$f$; Section~\ref{sec:sublinear} gives the changes needed under \textup{(E1)--(E3)}, the main one being a moment
bound for the in-degree of a vertex which replaces the uniform bound on the probability of an edge.
Section~\ref{sec:sim} compares the formulas with simulations.

\section{A priori bounds}\label{sec:apriori}

Throughout this section $f$ satisfies $0<\alpha\le f(k)\le\beta<\infty$ for all $k\ge0$, and $C=\beta/\alpha$.
The $m$ edges sent by a vertex $t$ are called the \emph{slots} of $t$ and are numbered $1,\dots,m$ in the order
in which their endpoints are chosen.

\begin{lemma}[Joint edge bound]\label{lem:edges}
Let $E$ be a finite multiset of pairs $\{i,j\}$ of distinct labels in $[n]$, with $e$ elements counted with
multiplicity. The probability that $G_n$ contains every pair of $E$ as an edge, with at least the prescribed
multiplicity, is at most
\[
   \prod_{\{i,j\}\in E}\frac{mC}{\max(i,j)-1}\;\le\;(2mC)^e\prod_{\{i,j\}\in E}\frac1{\max(i,j)} .
\]
\end{lemma}

\begin{proof}
At any slot of a vertex $t$ and for any history, a given vertex $i<t$ is chosen with conditional probability
$f(D_i)/S\le\beta/(\alpha(t-1))=C/(t-1)$, because the denominator is a sum of $t-1$ weights, each at least $\alpha$.
Assign to every element $\{i,j\}$ of $E$ a slot of its younger endpoint $\max(i,j)$, distinct elements receiving
distinct slots. Exposing the slots in chronological order and conditioning successively on the past, the
probability that every assigned slot chooses its prescribed endpoint is at most the product of $C/(\max(i,j)-1)$
over $E$. There are at most $m^e$ assignments, and the event in the statement is the union of the corresponding
events. The second inequality uses $t-1\ge t/2$ for $t\ge2$.
\end{proof}

\begin{lemma}\label{lem:schur}
The symmetric $n\times n$ matrix $H=(1/\max(i,j))_{i,j\le n}$ has operator norm at most $4$ on $\ell^2$.
\end{lemma}

\begin{proof}
Let $w_i=i^{-1/2}$. Then $\sum_{j\le i}H_{ij}w_j=i^{-1}\sum_{j\le i}j^{-1/2}\le 2i^{-1/2}$ and
$\sum_{j>i}H_{ij}w_j=\sum_{j>i}j^{-3/2}\le 2i^{-1/2}$, so $Hw\le 4w$ coordinatewise. For real $x$, the inequality
$2|x_ix_j|\le x_i^2w_j/w_i+x_j^2w_i/w_j$ and the symmetry of $H$ give
$|x^{\mathsf T}Hx|\le\sum_i x_i^2 (Hw)_i/w_i\le4\|x\|^2$.
\end{proof}

Let $U_n$ be uniform on $[n]$ and independent of $G_n$, and let $B_r(G_n,U_n)$ be the ball of radius $r$ around
$U_n$ for the graph distance that ignores orientation. Put $B=8mC$.

\begin{proposition}[Balls are bounded in mean and avoid old vertices]\label{prop:balls}
For all $n\ge1$, $r\ge0$ and $1\le K\le n$,
\[
  \E\,|B_r(G_n,U_n)|\le\sum_{l=0}^{r}B^l,\qquad
  \Prob\bigl(B_r(G_n,U_n)\cap[K]\neq\emptyset\bigr)\le\Bigl(\sum_{l=0}^{r}B^l\Bigr)\sqrt{K/n}.
\]
\end{proposition}

\begin{proof}
A vertex at distance $l$ from $U_n$ is the endpoint of a path $v_0=U_n,v_1,\dots,v_l$ with distinct vertices,
hence distinct edges. By Lemma~\ref{lem:edges}, the expected number of such paths ending in a set $A$ is at most
\[
  \frac{(2mC)^l}{n}\sum_{v_0,\dots,v_l}\prod_{k=1}^{l}\frac{1}{\max(v_{k-1},v_k)}\mathbf 1_{v_l\in A}
  =\frac{(2mC)^l}{n}\,\mathbf 1^{\mathsf T}H^l\mathbf 1_A\le\frac{(2mC)^l}{n}\,\sqrt n\,4^l\sqrt{|A|},
\]
by Lemma~\ref{lem:schur}. Take $A=[n]$ and $A=[K]$ and sum over $l\le r$.
\end{proof}

\begin{proposition}[No short cycles]\label{prop:cycles}
For every fixed $r$, the probability that $B_r(G_n,U_n)$ contains a cycle or a multiple edge tends to $0$.
\end{proposition}

\begin{proof}
Fix $\eps\in(0,1)$ and let $K=\lfloor\eps n\rfloor$. Suppose that the ball avoids $[K]$ and is not a simple tree.
Take a breadth-first spanning tree of the ball and an edge of the ball outside it (a second copy of a multiple edge
counts as such an edge). This edge together with the tree paths from its endpoints to the root forms a connected
rooted multigraph with $v\le 2r+1$ vertices and exactly $v$ edges, all of whose labels exceed $K$. There are
finitely many such rooted shapes. For a fixed shape there are at most $n^v$ labellings; the root label equals
$U_n$ with probability $1/n$; and by Lemma~\ref{lem:edges} the $v$ edges are present with probability at most
$(2mC/(\eps n))^v$. The union bound gives $O_{r,\eps}(1/n)$. With Proposition~\ref{prop:balls},
\[
  \limsup_{n\to\infty}\Prob\bigl(B_r(G_n,U_n)\text{ is not a simple tree}\bigr)\le\Bigl(\sum_{l\le r}B^l\Bigr)\sqrt\eps ,
\]
and $\eps$ is arbitrary.
\end{proof}

\begin{corollary}\label{cor:tight}
The laws of the rooted graphs $(G_n,U_n)$ are tight for the local topology. Every subsequential limit is almost
surely a locally finite simple tree in which every vertex has out-degree $m$; in particular it is infinite.
\end{corollary}

\begin{proof}
A vertex has at most $m$ outgoing edges, so there are finitely many rooted directed multigraphs with at most $M$
vertices that can occur as balls. Tightness follows from the first bound of Proposition~\ref{prop:balls} and
Markov's inequality, applied for each $r$ with a threshold $M_r$ chosen so that the exceptional probabilities are
summable. Limits are simple trees by Proposition~\ref{prop:cycles}. Every vertex other than vertex $1$ has
out-degree $m$, all edges leaving $B_r$ are seen in $B_{r+1}$, and $B_{r+1}$ contains vertex $1$ with probability
at most $(\sum_{l\le r+1}B^l)n^{-1/2}$; hence all out-degrees in a limit equal $m$. A finite tree with $v$
vertices has $v-1<mv$ edges, so the limit is infinite.
\end{proof}

\begin{lemma}[Degrees]\label{lem:degrees}
Let $D_i(n)$ be the in-degree of $i$ in $G_n$. Then $\E D_{U_n}(n)^2\le mC+2(mC)^2$, and for every $\delta>0$,
\[
   \Prob\Bigl(\max_{i\le n}D_i(n)\ge\delta n\Bigr)\le\frac{mC+2(mC)^2}{\delta^2n}.
\]
\end{lemma}

\begin{proof}
Let $L_i=mC\sum_{s=i}^{n-1}1/s$. Summing the single-slot bound $C/(t-1)$ over the slots of $t=i+1,\dots,n$ gives
$\E D_i(n)\le L_i$, and summing the product of two such bounds over ordered pairs of distinct slots gives
$\E[D_i(n)(D_i(n)-1)]\le L_i^2$. Exchanging the order of summation,
\[
  \sum_{i\le n}L_i=mC(n-1),\qquad \sum_{i\le n}L_i^2=(mC)^2\sum_{s,t<n}\frac{\min(s,t)}{st}\le 2(mC)^2(n-1).
\]
Hence $\sum_i\E D_i(n)^2\le n\,(mC+2(mC)^2)$, which gives the first claim, and the second follows from Markov's
inequality and a union bound.
\end{proof}

The second statement says that no vertex receives a positive fraction of the edges. The first gives uniform
integrability of the degree of the root, so the root of every subsequential limit has mean in-degree $m$.

\section{The normalisation and the degree law}\label{sec:normal}

In this section $0<\alpha\le f\le\beta<\infty$. Let $N_k(n)$ be the number of vertices of $G_n$ with in-degree $k$,
$y_n=(N_k(n)/n)_{k\ge0}$, and $S_n=\sum_{i\le n}f(D_i(n))=n\,s(y_n)$, where $s(y)=\sum_kf(k)y_k$.

\subsection{The scalar equation}

For $d>0$ let
\[
  R_k(d)=\prod_{j=0}^{k-1}\frac{mf(j)}{d+mf(j)}\quad(k\ge1),\qquad M(d)=\sum_{k\ge1}R_k(d),
\]
and let $p(d)=(p_k(d))_{k\ge0}$ be the law of $X(T)$, where $X$ is the pure birth process started at $0$ that
jumps from $k$ to $k+1$ at rate $(m/d)f(k)$, and $T$ is an independent exponential time of mean one.

\begin{lemma}\label{lem:scalar}
\textup{(i)} $\Prob(X(T)\ge k)=R_k(d)$, so $p(d)$ has mean $M(d)$, and
\begin{equation}\label{eq:prec}
  p_0(d)=\frac{d}{d+mf(0)},\qquad d\,p_k(d)=mf(k-1)p_{k-1}(d)-mf(k)p_k(d)\quad(k\ge1).
\end{equation}
\textup{(ii)} $d\,M(d)=m\sum_{k\ge0}f(k)p_k(d)$.
\textup{(iii)} There is exactly one $c_f>0$ with $M(c_f)=m$, and $c_f\in[\alpha,\beta]$.
\end{lemma}

\begin{proof}
(i) From state $j$ the process jumps before the killing time with probability $mf(j)/(d+mf(j))$, and
\eqref{eq:prec} is the same computation written for the point masses. (ii) $R_k(d)\le(m\beta/(d+m\beta))^k$, so all
sums converge absolutely; multiply \eqref{eq:prec} by $k$ and sum over $k\ge1$. (iii) Each $R_k$ is continuous and
strictly decreasing in $d$, the series converges uniformly on $[d_0,\infty)$ for every $d_0>0$,
$M(d)\le m\beta/d\to0$ as $d\to\infty$, and $M(d)\ge\sum_{k\le K}R_k(d)\to K$ as $d\to0$ for every $K$. At
$d=\alpha$ every factor is at least $m/(m+1)$, so $M(\alpha)\ge m$; at $d=\beta$ every factor is at most $m/(m+1)$,
so $M(\beta)\le m$.
\end{proof}

\begin{theorem}\label{thm:normal}
Almost surely $y_n\to p(c_f)$ in $\ell^1$ and $S_n/n\to c_f$.
\end{theorem}

The rest of the section proves Theorem~\ref{thm:normal}. No monotonicity of $f$ is used.

\subsection{Stochastic approximation in $\ell^1$}

Let $e_0$ be the unit mass at $0$ and let $Q$ be the bounded operator on $\ell^1$ given by
\[
  (Qy)_0=-f(0)y_0,\qquad (Qy)_k=f(k-1)y_{k-1}-f(k)y_k\quad(k\ge1),
\]
the generator of the pure birth process with rates $f(k)$ acting on measures; $\|Q\|\le2\beta$ and
$\sum_k(Qy)_k=0$. On probability vectors put
\[
   F(y)=e_0+\frac{m}{s(y)}\,Qy-y .
\]
Since $s\ge\alpha$ and $|s(y)-s(z)|\le\beta\|y-z\|_1$, $F$ is Lipschitz in $\ell^1$ on probability vectors.

At a slot of vertex $n+1$, if the current counts are $N$ and the current total weight is $S$, a vertex of
in-degree $k$ is chosen with probability $f(k)N_k/S$, so the conditional mean increment of the count vector is
$QN/S$. During the $m$ slots $N$ changes by at most $2m$ in $\ell^1$ and $S$ by at most $2m\beta$, while
$S\ge\alpha n$ and $\|QN\|_1\le2\beta n$; hence $QN/S$ differs from $Qy_n/s(y_n)$ by $O(1/n)$ in $\ell^1$,
uniformly in the configuration. After the $m$ slots the new vertex adds $e_0$. Therefore
\begin{equation}\label{eq:sa}
   y_{n+1}=y_n+\frac{1}{n+1}\bigl(F(y_n)+\eta_{n+1}+\rho_n\bigr),
\end{equation}
where $\E[\eta_{n+1}\mid\cF_n]=0$ coordinatewise, $\|\eta_{n+1}\|_1\le4m+2$ and $\|\rho_n\|_1\le K_0/n$ for a
constant $K_0$.

Every $y_n$ lies in
\[
   \cK=\Bigl\{y\ge0:\ \sum_ky_k=1,\ \sum_kk\,y_k\le m\Bigr\},
\]
because the total in-degree of $G_n$ is $m(n-1)$. The set $\cK$ is compact in $\ell^1$: the mass above $J$ is at
most $m/(J+1)$ uniformly on $\cK$, and the mean is lower semicontinuous. For $y\in\cK$ and $J\ge1$,
\begin{equation}\label{eq:Ftail}
   \|F(y)\|_1\le2+\frac{2m\beta}{\alpha},\qquad \sum_{k>J}|F(y)_k|\le\Bigl(1+\frac{2m\beta}{\alpha}\Bigr)\frac mJ .
\end{equation}

\begin{lemma}[Noise]\label{lem:noise}
Let $\tau_n=\sum_{j<n}1/(j+1)$. Almost surely, for every $T<\infty$,
\[
   \lim_{N\to\infty}\ \sup_{j\ge N:\ \tau_j-\tau_N\le T}\ \Bigl\|\sum_{n=N}^{j-1}\frac{\eta_{n+1}}{n+1}\Bigr\|_1=0 .
\]
\end{lemma}

\begin{proof}
For each coordinate $k$ the series $\sum_n\eta_{n+1,k}/(n+1)$ is a martingale with summable increment variances,
so it converges almost surely; work on the event where this holds for all $k$. The coordinates $k\le J$ of the
partial sums in the statement then tend to $0$ uniformly in $j$. For the coordinates $k>J$ use \eqref{eq:sa}
summed from $N$ to $j-1$: the noise sum equals $y_j-y_N$ minus the sums of the $F$ terms and of the $\rho$ terms.
The mass of $y_j$ and $y_N$ above $J$ is at most $m/(J+1)$ each, the $F$ terms contribute at most a constant times
$(T+1)/J$ by \eqref{eq:Ftail}, and the $\rho$ terms at most $K_0\sum_{n\ge N}1/(n(n+1))$. Hence the $\limsup$ in
$N$ of the supremum is at most a constant times $(T+1)/J$, and $J$ is arbitrary.
\end{proof}

The argument avoids any convergence theorem for $\ell^1$-valued martingales: the tails are controlled by the
recursion itself.

\subsection{The limiting flow}

Let $P_t=e^{tQ}$ be the semigroup of the pure birth process with rates $f(k)$; it is positive, preserves mass,
and $\|P_t\|=1$ on $\ell^1$. Let
\[
   A=Q+\frac1m\,e_0\,s(\cdot),\qquad w(t)=e^{tA}y,
\]
for a probability vector $y$. The operator $A$ is bounded and $e^{tA}$ is positive.

\begin{lemma}\label{lem:flowrep}
Let $|w|=\sum_kw_k$ and $z(t)=w(t)/|w(t)|$. Then $z'=\frac{s(z)}{m}F(z)$. Consequently the solution of
$\dot y=F(y)$ started at a probability vector $y$ is $\phi_\tau(y)=z(t(\tau))$, where
$\mathrm d\tau/\mathrm dt=s(z(t))/m\in[\alpha/m,\beta/m]$; it exists for all times, stays in the probability
vectors, and $\phi_\tau(\cK)\subset\cK$.
\end{lemma}

\begin{proof}
Since $Q$ preserves mass, $|w|'=s(w)/m$, and therefore
$z'=Az-\frac{s(z)}mz=Qz+\frac{s(z)}m(e_0-z)=\frac{s(z)}mF(z)$. Uniqueness for $\dot y=F(y)$ follows from the
Lipschitz property. For $y$ with finite mean, $\sum_kk\,(Qy)_k=s(y)$, so the mean $\mu(\tau)$ of $\phi_\tau(y)$
satisfies $\mu'=m-\mu$, and $\mu(0)\le m$ implies $\mu(\tau)\le m$. (The operator $Q$ is bounded for the weighted
norm $\sum_k(1+k)|y_k|$, so the mean stays finite and the computation is legitimate.)
\end{proof}

\begin{proposition}[Uniform attraction]\label{prop:attract}
$\sup_{y\in\cK}\|\phi_\tau(y)-p(c_f)\|_1\to0$ as $\tau\to\infty$.
\end{proposition}

\begin{proof}
Let $g_y(t)=s(P_ty)$ and $h(t)=s(P_te_0)$; both take values in $[\alpha,\beta]$ and are Lipschitz. Variation of
constants gives
\begin{equation}\label{eq:voc}
  w(t)=P_ty+\frac1m\int_0^tP_{t-u}e_0\,\sigma(u)\,\mathrm du,\qquad
  \sigma(t)=g_y(t)+\frac1m\int_0^th(t-u)\,\sigma(u)\,\mathrm du,
\end{equation}
with $\sigma(t)=s(w(t))$. The function $\gamma\mapsto\frac1m\int_0^\infty e^{-\gamma t}h(t)\,\mathrm dt$ is
continuous and strictly decreasing, and lies between $\alpha/(m\gamma)$ and $\beta/(m\gamma)$. Let $\gamma>0$ be
the value at which it equals one; then $\gamma\in[\alpha/m,\beta/m]$, and
$\kappa(t)=e^{-\gamma t}h(t)/m$ is a continuous, positive probability density on $[0,\infty)$ with finite mean
$\bar\mu$, bounded by $(\beta/m)e^{-\gamma t}$, hence directly Riemann integrable. Let
$v=\sum_{j\ge1}\kappa^{*j}$ be the renewal density; the series converges locally uniformly since
$\kappa^{*j}(t)\le(\beta/m)^jt^{j-1}/(j-1)!$. The key renewal theorem \cite[Chapter~XI]{Feller}, applied to the
function $\kappa$ itself, gives $v(t)\to1/\bar\mu$ as $t\to\infty$; in particular $v$ is bounded.

Multiplying the second equation of \eqref{eq:voc} by $e^{-\gamma t}$ gives a renewal equation for
$\tilde\sigma(t)=e^{-\gamma t}\sigma(t)$ with $\tilde g_y(t)=e^{-\gamma t}g_y(t)$, whose solution is
\[
   \tilde\sigma(t)=\tilde g_y(t)+\int_0^tv(t-u)\,\tilde g_y(u)\,\mathrm du .
\]
Since $\alpha e^{-\gamma u}\le\tilde g_y(u)\le\beta e^{-\gamma u}$ for every probability vector $y$, dominated
convergence gives
\[
   \tilde\sigma(t)\to C_y:=\frac1{\bar\mu}\int_0^\infty\tilde g_y(u)\,\mathrm du\in
   \Bigl[\frac{\alpha}{\gamma\bar\mu},\frac{\beta}{\gamma\bar\mu}\Bigr]
\]
uniformly in $y$. Multiplying the first equation of \eqref{eq:voc} by $e^{-\gamma t}$ and using
$\|P_ue_0\|_1=1$, we get $e^{-\gamma t}w(t)\to\frac{C_y}{m}\,\cR$ in $\ell^1$, uniformly in $y$, where
$\cR=\int_0^\infty e^{-\gamma u}P_ue_0\,\mathrm du$ has mass $1/\gamma$. Hence $z(t)\to\pi:=\gamma\cR$ uniformly
in $y$, and by Lemma~\ref{lem:flowrep} (the clock $\tau(t)$ runs at speed at least $\alpha/m$)
$\phi_\tau(y)\to\pi$ uniformly in $y$.

The vector $\pi$ is the law of the pure birth process with rates $f(k)$ at an independent exponential time of
rate $\gamma$, that is, $\pi=p(m\gamma)$. Its mean is
$\int_0^\infty\gamma e^{-\gamma t}\int_0^th(u)\,\mathrm du\,\mathrm dt=\int_0^\infty e^{-\gamma u}h(u)\,\mathrm du=m$
by the choice of $\gamma$. Thus $M(m\gamma)=m$, and Lemma~\ref{lem:scalar}(iii) gives $m\gamma=c_f$.
\end{proof}

\begin{proof}[Proof of Theorem~\ref{thm:normal}]
Interpolate $y_n$ linearly at the times $\tau_n$ and call the path $y(\cdot)$. By \eqref{eq:sa},
Lemma~\ref{lem:noise}, the bound $\|\rho_n\|_1\le K_0/n$, the Lipschitz property of $F$ and Gronwall's inequality,
almost surely for every $T$
\[
   \sup_{0\le u\le T}\bigl\|y(\tau+u)-\phi_u(y(\tau))\bigr\|_1\to0\qquad(\tau\to\infty);
\]
the replacement of the step sums of $F(y_n)$ by an integral costs $O(\sum_{n\ge N}(n+1)^{-2})$ because
$\|y_{n+1}-y_n\|_1=O(1/n)$. Given $\eps>0$ choose $T$ with $\sup_{y\in\cK}\|\phi_T(y)-p(c_f)\|_1<\eps$
(Proposition~\ref{prop:attract}); the interpolated path stays in the convex set $\cK$, so
$\limsup_{\tau}\|y(\tau)-p(c_f)\|_1\le\eps$ almost surely. Finally $S_n/n=s(y_n)\to s(p(c_f))=c_f$ by
Lemma~\ref{lem:scalar}(ii), since $s$ is continuous on $\ell^1$.
\end{proof}

\begin{example}[Two levels]\label{ex:two}
Let $f(0)=a$ and $f(k)=b$ for $k\ge1$. Then $S_n=bn+(a-b)N_0(n)$, and $z=p_0(c_f)$ is the root in $(0,1)$ of
\[
   (a-b)z^2+(2b+ma-a)z-b=0,\qquad c_f=b+(a-b)z ,
\]
with $p_1=\lambda az/(1+\lambda b)$ and $p_k=p_1\,(\lambda b/(1+\lambda b))^{k-1}$, $\lambda=m/c_f$. For $m=2$,
$a=1$, $b=2$ this gives $z=(5-\sqrt{17})/2$ and $c_f=(\sqrt{17}-1)/2$. For $a=b$ the law is geometric with ratio
$m/(m+1)$, the degree law of uniform attachment.
\end{example}

\section{The local limit}\label{sec:local}

In this section $0<\alpha\le f\le\beta<\infty$, $c=c_f$ and $\lambda=m/c$. By Theorem~\ref{thm:normal},
$S_t/t\to c$ almost surely; since $\alpha\le S_t/t\le\beta$, for every $\eps>0$
\begin{equation}\label{eq:tailsup}
   \E\sup_{t\ge\eps n}\bigl|S_t/t-c\bigr|\to0\qquad(n\to\infty).
\end{equation}

\subsection{Ball shapes and the limiting probabilities}

Fix $r\ge1$. A \emph{ball shape} is a finite rooted directed simple tree $H$ with vertex set $V$ and root $o$ in
which every vertex is within distance $r$ of $o$ and every vertex at distance $r$ is a leaf. Let $I$ be the set
of vertices at distance less than $r$ (the interior), and for $j\in V$ let $q_j$ and $k_j$ be the out-degree and
the in-degree of $j$ in $H$. We write $j\to i$ for an edge sent by $j$ and received by $i$. The shape is
\emph{admissible} if $q_j=m$ for every $j\in I$ and $q_j\le m$ for every $j$. Let $A(H)$ be the number of
automorphisms of $H$ that fix the root and preserve orientations, and $(m)_q=m!/(m-q)!$.

An \emph{age vector} is $x\in(0,1)^V$ with distinct coordinates; it is \emph{compatible} if $x_j>x_i$ for every
edge $j\to i$. Given a compatible age vector, let $(Y_i)_{i\in V}$ be independent processes, where $Y_i$ is
defined on $[x_i,1]$, starts at $0$, jumps by $+1$ at rate $\lambda f(Y_i(t))/t$ at time $t$, and in addition
jumps by $+1$ at each of the times $x_j$, $j\to i$. Define
\begin{equation}\label{eq:Phi}
  \Phi_H(x)=\prod_{j\in V}(m)_{q_j}\ \prod_{j\to i}\frac{1}{c\,x_j}\ \,
     \E\Bigl[\,\prod_{j\to i}f\bigl(Y_i(x_j-)\bigr)\prod_{i\in I}\mathbf 1_{\{Y_i(1)=k_i\}}\Bigr]
\end{equation}
for compatible $x$ and $\Phi_H(x)=0$ otherwise, and
\begin{equation}\label{eq:PH}
   \mathbf P_H=\frac{1}{A(H)}\int_{(0,1)^V}\Phi_H(x)\,\mathrm dx .
\end{equation}

\begin{theorem}\label{thm:localbdd}
For every $r\ge1$ and every admissible ball shape $H$ of radius $r$,
\[
   \frac1n\,\#\bigl\{u\in[n]:\ B_r(G_n,u)\cong H\bigr\}\ \longrightarrow\ \mathbf P_H\qquad\text{in probability,}
\]
and the proportion of $u$ whose ball $B_r(G_n,u)$ is not an admissible ball shape tends to $0$ in probability.
Consequently $G_n$ converges locally in probability to the random rooted tree whose balls have the laws
$(\mathbf P_H)$; this tree is infinite, simple, and all its vertices have out-degree $m$.
\end{theorem}

The expectation in \eqref{eq:Phi} is a product of explicit factors.

\begin{proposition}\label{prop:factors}
Let $X$ be the pure birth process started at $0$ with rates $\lambda f(k)$, and $h_\lambda(u)=\E f(X(u))$. For
compatible $x$ the expectation in \eqref{eq:Phi} equals $\prod_{i\in V}F_i$, where:
for $i\in I$ with children of ages $t_1<\dots<t_k$ $(k=k_i)$, and $t_0=x_i$, $t_{k+1}=1$,
\[
   F_i=\prod_{\ell=0}^{k-1}f(\ell)\ \prod_{\ell=0}^{k}\Bigl(\frac{t_\ell}{t_{\ell+1}}\Bigr)^{\lambda f(\ell)};
\]
for $i\notin I$, $F_i=1$ if $k_i=0$ and $F_i=h_\lambda(\log(t_1/x_i))$ if $k_i=1$, where $t_1$ is the age of the
child of $i$.
\end{proposition}

\begin{proof}
The processes are independent, so the expectation factorises. For $i\in I$ the event $Y_i(1)=k_i$ says that
$Y_i$ makes no jump other than the $k_i$ forced ones; on this event $Y_i=\ell$ on $[t_\ell,t_{\ell+1})$, the
probability of no other jump on that interval is $\exp(-\lambda f(\ell)\int_{t_\ell}^{t_{\ell+1}}\mathrm dt/t)$,
and the weight at the $\ell$th forced jump is $f(\ell-1)$. A vertex $i\notin I$ is a leaf of $H$ with no
constraint at time $1$; if its neighbour is its child of age $t_1$, then $Y_i(t_1-)$ has the law of
$X(\log(t_1/x_i))$.
\end{proof}

For instance, if $H$ is the star in which the root has no child, the root contributes $x_o^{\lambda f(0)}$, and
each of its $m$ parents, of age $y<x_o$, contributes $h_\lambda(\log(x_o/y))/(c\,x_o)$, which integrates to
$\int_0^\infty e^{-u}h_\lambda(u)\,\mathrm du/c=1$ over $y\in(0,x_o)$ by Lemma~\ref{lem:scalar}(ii); since
$(m)_m=A(H)=m!$, this gives $\mathbf P_H=\int_0^1x^{\lambda f(0)}\,\mathrm dx=p_0(c)$, in agreement with
Theorem~\ref{thm:normal}. The same computation shows that the age $y$ of a parent of a vertex of age $x$ has
density $h_\lambda(\log(x/y))/(cx)$ on $(0,x)$, which is not uniform unless $f$ is constant.

\subsection{Forced slots}

Fix an admissible $H$ with $e=|V|-1$ edges, an injective labelling $\iota:V\to[n]$ with $\iota(j)>\iota(i)$ for
every edge $j\to i$, and a \emph{slot assignment} $\sigma$, which gives to every edge $j\to i$ a slot of
$\iota(j)$, different edges receiving different slots. There are $\prod_j(m)_{q_j}$ slot assignments. Let
$E_{\iota,\sigma}$ be the event that for every edge $j\to i$ the assigned slot of $\iota(j)$ chooses
$\iota(i)$, and let $D_\iota=\{D_{\iota(i)}(n)=k_i\ \text{for all }i\in I\}$.

The \emph{forced process} is the graph process in which every assigned slot chooses its prescribed endpoint with
probability one and all other slots follow the original rule. Write $\Prob^{\sigma}$, $\E^\sigma$ for its law.
Multiplying the transition probabilities along a history gives, for every nonnegative function $\Psi$ of the
history up to time $n$,
\begin{equation}\label{eq:com}
   \E\bigl[\Psi\,\mathbf 1_{E_{\iota,\sigma}}\bigr]
   =\E^\sigma\Bigl[\Psi\prod_{j\to i}\frac{f(D^-_{j\to i})}{S^-_{j\to i}}\Bigr],
\end{equation}
where $D^-_{j\to i}$ is the in-degree of $\iota(i)$ and $S^-_{j\to i}$ the total weight of the eligible vertices
immediately before the slot assigned to $j\to i$.

\begin{lemma}[The forced process has the same normalisation]\label{lem:forcedS}
Suppose that all labels $\iota(j)$ are at least $\eps n$. Under $\Prob^\sigma$,
$\sup_{\eps n\le t\le n}|S_t/t-c|\to0$ in $L^1$ as $n\to\infty$, uniformly in $\iota$ and $\sigma$.
\end{lemma}

\begin{proof}
Couple the forced process with the original one: they coincide until the first assigned slot, and afterwards,
at each slot that is not assigned, the two endpoints are chosen by a maximal coupling of the two conditional
laws. Let $\bar K$ be the number of labels at which the two in-degree sequences have differed at some time. If
the two weight vectors at a slot of vertex $t+1$ are $w$ and $w'$ with totals $W,W'\ge\alpha t$, then
\[
   \bigl\|w/W-w'/W'\bigr\|_1\le\frac{2\|w-w'\|_1}{W}\le\frac{2\beta\bar K}{\alpha t},
\]
so the two choices differ with probability at most $C\bar K/t$, in which case $\bar K$ increases by at most $2$;
at an assigned slot $\bar K$ increases by at most $2$. Hence
$\E\bar K\le 2e\exp(2mC\sum_{t\ge\eps n/2}^{n}1/t)\le K_\eps$, a constant that does not depend on $n$. Since the
two values of $S_t$ differ by at most $\beta\bar K$, the claim follows from \eqref{eq:tailsup}.
\end{proof}

\begin{lemma}[Tagged degrees]\label{lem:tagged}
Let $\iota=\iota_n$ and $\sigma$ be as above, and suppose that $\iota_n(j)/n\to x_j$ for all $j$, where $x$ is a
compatible age vector with all $x_j\ge\eps$. Under $\Prob^\sigma$, the vector formed by the in-degrees
$D^-_{j\to i}$ $(j\to i)$ and $D_{\iota(i)}(n)$ $(i\in V)$ converges in total variation to the vector formed by
$Y_i(x_j-)$ and $Y_i(1)$.
\end{lemma}

\begin{proof}
The vertices $\iota(i)$, $i\in V$, are called tagged. Consider a slot of a vertex $t+1$ that is not assigned.
Given the history, the tagged vertex $\iota(i)\le t$ is chosen with probability $f(D_{\iota(i)})/S$, where $S$
is the current total weight, and at most one tagged vertex is chosen. Compare with independent Bernoulli
variables of parameters $f(D_{\iota(i)})/(ct)$: the total variation distance between the two laws of the
increment of the tagged vector is at most
\[
   K_1\,\frac{|S/t-c|}{t}+\frac{K_1}{t^2},
\]
with $K_1$ depending on $|V|,m,\alpha,\beta$ only, because $S\ge\alpha t$, $c\ge\alpha$ and the probability
that two of the Bernoulli variables equal one is $O(t^{-2})$. At an assigned slot the prescribed tagged vertex
gains one and the others gain nothing, in both systems. Using maximal couplings slot by slot, the tagged vector
of the forced process coincides up to time $n$ with a system of independent chains, one for each tagged vertex,
which start at $0$ at the birth of the vertex, jump at each unassigned slot of a vertex $t+1$ with probability
$f(\text{state})/(ct)$, and jump at their assigned slots, except on an event of probability at most
\[
   K_1m\log(1/\eps)\ \E^\sigma\!\sup_{t\ge\eps n}|S_t/t-c|+O\bigl(1/(\eps n)\bigr),
\]
which tends to $0$ by Lemma~\ref{lem:forcedS} (within a batch $S$ differs from $S_t$ by at most $2m\beta$).

It remains to compare each independent chain with $Y_i$. Over the slots of vertex $t+1$ the chain in state $k$
jumps exactly once with probability $mf(k)/(ct)+O(t^{-2})$ and more than once with probability $O(t^{-2})$;
the process $Y_i$ in state $k$ has the same two estimates on the age interval $[t/n,(t+1)/n]$, since
$\int\lambda f(k)\,\mathrm ds/s=\lambda f(k)\log(1+1/t)$. The total variation distance between the two paths,
observed at the ages $t/n$, is therefore $O(\sum_{t\ge\eps n}t^{-2})=O(1/(\eps n))$, apart from the effect of the
forced times, which are $\iota_n(j)/n$ for the chain and $x_j$ for $Y_i$: moving a forced time by $\delta$
changes the law of the observed vector by at most the probability of an unforced jump in a window of length
$\delta$, which is $O(\delta/\eps)$. Finally the in-degree of a target immediately before its assigned slot
differs from its in-degree at the beginning of the batch with probability $O(1/n)$.
\end{proof}

\begin{proposition}\label{prop:embed}
Under the assumptions of Lemma~\ref{lem:tagged},
\[
   n^{e}\,\Prob\bigl(E_{\iota_n,\sigma}\cap D_{\iota_n}\bigr)\ \longrightarrow\
   \prod_{j\to i}\frac1{c\,x_j}\ \E\Bigl[\prod_{j\to i}f\bigl(Y_i(x_j-)\bigr)\prod_{i\in I}\mathbf 1_{\{Y_i(1)=k_i\}}\Bigr],
\]
and the left-hand side is bounded by $(4C/\eps)^e$ for all labellings with labels at least $\eps n$ and $n$
large. The convergence is uniform on every compact set of compatible age vectors.
\end{proposition}

\begin{proof}
Apply \eqref{eq:com} with $\Psi=\mathbf 1_{D_\iota}$. Each factor $nf(D^-)/S^-$ is at most
$n\beta/(\alpha(\eps n-1))$, and $S^-_{j\to i}/n\to c\,x_j$ in $L^1$ under $\Prob^\sigma$ by
Lemma~\ref{lem:forcedS}. The numerators and the indicator are functions of the vector of
Lemma~\ref{lem:tagged}. The limit is continuous in $x$ on the set of compatible age vectors, and convergence
along every sequence $\iota_n(\cdot)/n\to x$ gives the uniformity.
\end{proof}

\subsection{Proof of Theorem~\ref{thm:localbdd}}

\emph{Counting.} Let $\mathcal A_\iota$ be the event that $\iota$ is an isomorphism from $H$ onto
$B_r(G_n,\iota(o))$. A ball isomorphic to $H$ has exactly $A(H)$ such isomorphisms, so
\begin{equation}\label{eq:count}
   \Prob\bigl(B_r(G_n,U_n)\cong H\bigr)=\frac{1}{n\,A(H)}\sum_\iota\Prob(\mathcal A_\iota).
\end{equation}
We compare $\mathbf 1_{\mathcal A_\iota}$ with $\sum_\sigma\mathbf 1_{E_{\iota,\sigma}\cap D_\iota}$. Every edge
of $H$ has an endpoint in $I$. On $D_\iota$ no edge $j\to i$ with $i\in I$ can be doubled, because the in-degree
of $\iota(i)$ is exactly the number of its children; and an edge $j\to i$ with $j\in I$ cannot be doubled on
$E_{\iota,\sigma}$, because all $m$ slots of $\iota(j)$ are assigned to distinct endpoints. Hence at most one
$\sigma$ contributes, and on $\mathcal A_\iota$ exactly one does. Conversely, on $E_{\iota,\sigma}\cap D_\iota$
all edges of $G_n$ at the vertices $\iota(i)$, $i\in I$, are images of edges of $H$; by induction on the
distance, the vertices at distance $d\le r$ from $\iota(o)$ in $G_n$ are the images of the vertices at distance
$d$ in $H$. Thus $\mathcal A_\iota$ can fail only if $G_n$ has an edge between the images of two vertices at
distance $r$. If all labels are at least $\eps n$, Lemma~\ref{lem:edges} bounds the probability of
$E_{\iota,\sigma}$ together with one such extra edge by $(4mC/(\eps n))^{e+1}$. There are at most
$n^{|V|}=n^{e+1}$ labellings, so
\begin{equation}\label{eq:count2}
   \Prob\bigl(B_r(G_n,U_n)\cong H,\ B_r(G_n,U_n)\subset[\eps n,n]\bigr)
   =\frac{1}{A(H)}\ \frac{1}{n^{|V|}}\sum_{\iota\ge\eps n}\ \sum_\sigma n^{e}\,\Prob(E_{\iota,\sigma}\cap D_\iota)+O_\eps(1/n).
\end{equation}

\emph{Riemann sums.} By Proposition~\ref{prop:embed} the summands are bounded, and they converge uniformly on
compact sets of compatible age vectors; the limit does not depend on $\sigma$. Labellings for which two ages are
within $\delta$ form a proportion $O(\delta+1/n)$ of all labellings. Letting $n\to\infty$ and then $\delta\to0$,
the right-hand side of \eqref{eq:count2} converges to $A(H)^{-1}\int_{[\eps,1]^V}\Phi_H(x)\,\mathrm dx$.

\emph{The cutoff.} By Proposition~\ref{prop:balls} the left-hand side of \eqref{eq:count2} differs from
$\Prob(B_r(G_n,U_n)\cong H)$ by at most $(\sum_{l\le r}B^l)\sqrt\eps$. The integrals over $[\eps,1]^V$ increase as
$\eps$ decreases and are bounded by one, so $\Phi_H$ is integrable and
$\Prob(B_r(G_n,U_n)\cong H)\to\mathbf P_H$. Balls that are not admissible ball shapes have vanishing probability
by Proposition~\ref{prop:cycles} and Corollary~\ref{cor:tight}, and tightness shows that $(\mathbf P_H)_H$ is a
probability distribution for each $r$.

\emph{Two roots.} Let $U_n,U_n'$ be independent uniform vertices. Their balls of radius $r$ intersect with
probability at most $\E|B_{2r}(G_n,U_n)|/n=O(1/n)$. For two admissible shapes $H_1,H_2$ repeat the argument with
the disjoint union $H_1\sqcup H_2$, which has $|V_1|+|V_2|-2$ edges, and two root factors $1/n^2$. An edge
between the image of $H_1$ and the image of $H_2$ can only join vertices at distance $r$ from the two roots and
does not change either ball, so the counting step is unchanged. No forced jump links the two components, the
processes $Y_i$ are independent, and the limiting density is $\Phi_{H_1}\Phi_{H_2}$. Therefore
\[
   \Prob\bigl(B_r(G_n,U_n)\cong H_1,\ B_r(G_n,U_n')\cong H_2\bigr)\to\mathbf P_{H_1}\mathbf P_{H_2}.
\]
With $H_1=H_2=H$ this says that the second moment of $n^{-1}\#\{u:B_r(G_n,u)\cong H\}$ converges to
$\mathbf P_H^2$, and the first moment converges to $\mathbf P_H$; hence the convergence in probability. Since
balls with more than $M$ vertices have expected proportion at most $(\sum_{l\le r}B^l)/M$, the convergence for
each shape gives convergence in probability of the empirical average of every bounded function of the ball of
radius $r$, which is local convergence in probability. The properties of the limit were proved in
Corollary~\ref{cor:tight}. \qed

\section{Unbounded sublinear weights}\label{sec:sublinear}

We now assume only \textup{(E1)--(E3)}: $f\ge\alpha>0$, $|f(k+1)-f(k)|\le L$ and $f(k)/k\to0$. A bounded $f$ with
$\inf f>0$ satisfies them, with $L=\beta-\alpha$. By \textup{(E3)}, for every $\eps>0$ there is $a_\eps\ge1$ with
\begin{equation}\label{eq:lin}
    f(k)\le\eps\,(k+a_\eps)\qquad\text{for all }k\ge0 .
\end{equation}
The plan of the proof is the same, and we describe the changes section by section. The bounds $f\le\beta$ and
$C=\beta/\alpha$ were used in three places: the single-slot bound $C/(t-1)$ behind
Section~\ref{sec:apriori}, the boundedness of the operator $Q$ in Section~\ref{sec:normal}, and the couplings of
Section~\ref{sec:local}.

\subsection{Moments of a degree and the edge bound}

Let $E$ be a multiset of pairs of labels with a slot assignment as in the proof of Lemma~\ref{lem:edges}. The
forced process, in which each assigned slot chooses its prescribed endpoint, and the identity \eqref{eq:com} are
defined as in Section~\ref{sec:local}; the identity holds for every positive $f$.

\begin{lemma}\label{lem:moments}
Let $\theta\in(0,1/2)$ and let $p\ge1$, $q\ge0$ be integers. There are $a\ge1$ and $K<\infty$ such that, in the
original process and in every forced process with at most $q$ assigned slots, for all $i<j$,
\[
    \E^\sigma\bigl[(D_i+a)^p\bigr]\le K\,(j/i)^{p\theta},
\]
where $D_i$ is the in-degree of $i$ at any moment up to the last slot of vertex $j$.
\end{lemma}

\begin{proof}
Let $\eps=\alpha\theta/(2m)$ and choose $a\ge a_\eps$ with $(1+1/a)^{p-1}\le2$. Put $Z=D_i+a$. At a slot of a vertex
$t+1$ that is not assigned, $i$ is chosen with conditional probability at most $f(D_i)/(\alpha t)\le\eps Z/(\alpha t)$,
and $(Z+1)^p-Z^p\le p\,(1+1/a)^{p-1}Z^{p-1}$. Hence the conditional expectation of $Z^p$ after the slot is at most
$(1+2p\eps/(\alpha t))Z^p$, and over the $m$ slots of the vertex the factor is at most
$\exp(2pm\eps/(\alpha t))=\exp(p\theta/t)$. An assigned slot multiplies $Z^p$ by at most $(1+1/a)^p$. Since
$\sum_{s=i}^{j-1}1/s\le1+\log(j/i)$, the claim holds with $K=a^p(1+1/a)^{pq}e^{p\theta}$.
\end{proof}

\begin{lemma}[Joint edge bound]\label{lem:edges2}
Let $\theta\in(0,1/2)$. For every $e$ there is $K_e$ such that, for every multiset $E$ of $e$ pairs of distinct
labels, the probability that $G_n$ contains every pair of $E$ with at least the prescribed multiplicity is at most
\[
    K_e\prod_{\{i,j\}\in E}\max(i,j)^{\theta-1}\min(i,j)^{-\theta}.
\]
\end{lemma}

\begin{proof}
Fix a slot assignment, and let $\eps$, $a$ and $K$ be as in Lemma~\ref{lem:moments} and its proof for $p=q=e$,
so that $f(k)\le\eps(k+a)$. By \eqref{eq:com} with $\Psi=1$, the bound $S^-\ge\alpha(j-1)\ge\alpha j/2$,
H\"older's inequality with exponent $e$, and Lemma~\ref{lem:moments}, the
probability that all assigned slots choose their prescribed endpoints is at most
\[
    \prod_{j\to i}\frac{2\eps}{\alpha j}\Bigl(\E^\sigma\bigl[(D^-_{j\to i}+a)^e\bigr]\Bigr)^{1/e}
    \le\prod_{j\to i}\frac{2\eps K^{1/e}}{\alpha}\,j^{\theta-1}i^{-\theta},
\]
where $j\to i$ runs over the pairs of $E$ with $j>i$. There are at most $m^e$ slot assignments.
\end{proof}

Let $H_\theta(i,j)=\max(i,j)^{\theta-1}\min(i,j)^{-\theta}$. With $w_i=i^{-1/2}$,
\[
   \sum_{j\le i}H_\theta(i,j)w_j=i^{\theta-1}\sum_{j\le i}j^{-\theta-1/2}\le\frac{i^{-1/2}}{1/2-\theta},\qquad
   \sum_{j>i}H_\theta(i,j)w_j=i^{-\theta}\sum_{j>i}j^{\theta-3/2}\le\frac{i^{-1/2}}{1/2-\theta},
\]
so the argument of Lemma~\ref{lem:schur} gives $\|H_\theta\|\le2/(1/2-\theta)$. Take $\theta=1/4$. The proofs of
Section~\ref{sec:apriori} then go through with Lemma~\ref{lem:edges2} in place of Lemma~\ref{lem:edges}: for each
$r$ there is $K_r$ with
\[
    \E|B_r(G_n,U_n)|\le K_r,\qquad \Prob\bigl(B_r(G_n,U_n)\cap[K]\ne\emptyset\bigr)\le K_r\sqrt{K/n};
\]
if all labels exceed $\eps n$, every factor in Lemma~\ref{lem:edges2} is at most $1/(\eps n)$, which is what the
proof of Proposition~\ref{prop:cycles} uses; and Corollary~\ref{cor:tight} is unchanged. Lemma~\ref{lem:moments}
with $p=2$, $q=0$ gives $\E D_i(n)^2\le K(n/i)^{1/2}$, hence $\E D_{U_n}(n)^2\le2K$ and
$\Prob(\max_iD_i(n)\ge\delta n)\le2K/(\delta^2n)$, which replaces Lemma~\ref{lem:degrees}.

\subsection{The normalisation}

The operator $Q$ is no longer bounded on $\ell^1$, and $F$ is not Lipschitz. On the compact set $\cK$, however,
\eqref{eq:lin} gives for $y,z\in\cK$ and $J\ge1$
\begin{equation}\label{eq:stail}
\begin{gathered}
   \alpha\le s(y)\le\eps(m+a_\eps),\qquad \sum_{k\ge J}f(k)\,y_k\le m\sup_{k\ge J}\frac{f(k)}k,\\
   \sum_kf(k)|y_k-z_k|\le\max_{k<J}f(k)\,\|y-z\|_1+2m\sup_{k\ge J}\frac{f(k)}k .
\end{gathered}
\end{equation}
Since $\|Qy\|_1\le2\sum_kf(k)|y_k|$, the maps $s$, $Q$ and $F$ are uniformly continuous on $\cK$ for the
$\ell^1$ norm, $\|F(y)\|_1\le2+2m$, and $\sup_{y\in\cK}\sum_{k>J}|F(y)_k|\to0$ as $J\to\infty$.

\emph{The recursion.} Let $b_n=\max_{k\le mn}f(k)$, so that $b_n/n\to0$. During the slots of vertex $n+1$ the
count vector $N$ changes by at most $2m$ in $\ell^1$, so $QN$ changes by at most $4mb_n$, and $S$ changes by at
most $mL$ by \textup{(E2)}, while $S\ge\alpha n$ and $\|QN\|_1\le2S$. Hence \eqref{eq:sa} holds with the same
$\eta_{n+1}$ and with $\|\rho_n\|_1\le K_0(b_n+1)/n\to0$. In the proof of Lemma~\ref{lem:noise} the
$\rho$ terms now contribute at most $(T+1)\sup_{n\ge N}\|\rho_n\|_1$, and the $F$ terms are controlled by the
uniform tail bound above; the conclusion of the lemma holds.

\emph{The semigroup.} The pure birth process $X$ with rates $f(k)$ does not explode, because by \eqref{eq:lin}
$\E_k(X_t+a_\eps)\le(k+a_\eps)e^{\eps t}$. Its semigroup $P_t$ acts on $\ell^1$ as a strongly continuous
semigroup of positive contractions that preserve mass, and it is bounded and strongly continuous for the norm
$\|y\|_{(1)}=\sum_k(1+k)|y_k|$. The operator $Q$ is bounded from this weighted space to $\ell^1$. For a finitely
supported $y$, $P_ty-y=\int_0^tP_uQy\,\mathrm du$; by density this extends to all $y$ with
$\|y\|_{(1)}<\infty$, so such $y$ belong to the domain of the generator of $(P_t)$ on $\ell^1$, and the generator
acts on them as $Q$. For $y\in\cK$ put $g_y(t)=s(P_ty)$ and $h(t)=s(P_te_0)$; then, for every $\eps>0$,
\begin{equation}\label{eq:gh}
   \alpha\le g_y(t)\le\eps(m+a_\eps)e^{\eps t},\qquad \alpha\le h(t)\le\eps a_\eps e^{\eps t},
   \qquad \E_0X_t=\int_0^th(u)\,\mathrm du .
\end{equation}

\begin{lemma}\label{lem:flow2}
Let $y:[0,T]\to\cK$ be a solution of $y(\tau)=y(0)+\int_0^\tau F(y(u))\,\mathrm du$. Define the clock $t(\tau)$ by
$\mathrm dt/\mathrm d\tau=m/s(y(\tau))$, $t(0)=0$, let $z(t)=y(\tau(t))$ and
$w(t)=z(t)\exp\bigl(\frac1m\int_0^ts(z(u))\,\mathrm du\bigr)$. Then $w$ and $\sigma=s(w)$ satisfy \eqref{eq:voc}
with initial vector $y(0)$. In particular $y$ is determined by $y(0)$.
\end{lemma}

\begin{proof}
The integrand is continuous, so $y$ is continuously differentiable in $\ell^1$, and so are $z$ and $w$. As in
Lemma~\ref{lem:flowrep}, $w'=Qw+\frac1me_0\,s(w)$, where $w(t)$ has finite mean and therefore lies in the domain of
the generator. For fixed $t$ the function $u\mapsto P_{t-u}w(u)$ is then differentiable on $[0,t]$ with derivative
$P_{t-u}(w'(u)-Qw(u))=\frac1mP_{t-u}e_0\,\sigma(u)$, which gives the first equation of \eqref{eq:voc}. Applying the
positive linear functional $s$ to it gives the second. The functions $g_y$ and $h$ are locally bounded by
\eqref{eq:gh} and $\sigma$ is continuous, so $\sigma$ is the unique locally bounded solution of this Volterra
equation; it determines $w$, then $z=w/|w|$, then the clock.
\end{proof}

\emph{Uniform attraction.} By \eqref{eq:gh} the Laplace transform of $h$ is finite on $(0,\infty)$, continuous
and strictly decreasing, at least $\alpha/\gamma$ at $\gamma$, and tends to $0$ at infinity. Let $\gamma>0$ be the
point where it equals $m$ and $\kappa(t)=e^{-\gamma t}h(t)/m$. Taking $\eps<\gamma$ in \eqref{eq:gh}, $\kappa$ is
continuous and bounded by a multiple of $e^{-(\gamma-\eps)t}$, hence directly Riemann integrable, and
$\tilde g_y(u)=e^{-\gamma u}g_y(u)$ lies between $\alpha e^{-\gamma u}$ and $\eps(m+a_\eps)e^{-(\gamma-\eps)u}$ for
all $y\in\cK$. The proof of Proposition~\ref{prop:attract} now applies word for word to the functions $w$ of
Lemma~\ref{lem:flow2}: $z(t)\to\pi=\gamma\int_0^\infty e^{-\gamma u}P_ue_0\,\mathrm du$ in $\ell^1$ as
$t\to\infty$, uniformly in $y(0)\in\cK$, and the clock satisfies $t\ge m\tau/(\eps(m+a_\eps))$. The vector $\pi$ is
$p(m\gamma)$; by the last identity of \eqref{eq:gh} its mean is $\int_0^\infty e^{-\gamma u}h(u)\,\mathrm du=m$,
and $s(\pi)=\gamma\int_0^\infty e^{-\gamma u}h(u)\,\mathrm du=m\gamma$. The function $M$ of
Section~\ref{sec:normal} is finite, since $M(d)=\int_0^\infty e^{-t}\,\E_0X_{mt/d}\,\mathrm dt$ and
$\E_0X_t\le a_\eps e^{\eps t}$ for every $\eps>0$, and it is strictly decreasing, so $c_f=m\gamma$
is the only solution of $M(d)=m$. This proves Lemma~\ref{lem:scalar}(iii) in the present setting, without the
bounds $\alpha\le c_f\le\beta$; one has $c_f\ge\alpha$.

\emph{Proof of Theorem~\ref{thm:normal}.} Suppose that, with positive probability,
$\|y(\tau_k)-\pi\|_1>2\delta$ along a sequence $\tau_k\to\infty$. Fix $T$ such that every solution
$y:[0,T]\to\cK$ of the integral equation of Lemma~\ref{lem:flow2} has $\|y(T)-\pi\|_1\le\delta$. The paths
$u\mapsto y(\tau_k-T+u)$, $u\in[0,T]$, take values in the compact set $\cK$ and are uniformly Lipschitz, so a
subsequence converges uniformly to some $\bar y$. By \eqref{eq:sa}, Lemma~\ref{lem:noise}, $\|\rho_n\|_1\to0$ and
the uniform continuity of $F$ on $\cK$, the limit satisfies $\bar y(u)=\bar y(0)+\int_0^uF(\bar y(v))\,\mathrm dv$.
Hence $\|\bar y(T)-\pi\|_1\le\delta$, a contradiction. Thus $y_n\to\pi=p(c_f)$ almost surely in $\ell^1$, and
$S_n/n=s(y_n)\to s(\pi)=c_f$ because $s$ is continuous on $\cK$.

\subsection{The local limit}

Since $s$ is bounded on $\cK$, \eqref{eq:tailsup} holds. In Lemma~\ref{lem:forcedS} replace $\bar K$ by a bound
for the $\ell^1$ distance $\bar d$ between the two in-degree sequences: by \textup{(E2)} the two weight vectors
differ by at most $L\bar d$ in $\ell^1$, the two choices at a slot of vertex $t+1$ differ with probability at most
$L\bar d/(\alpha t)$, a disagreement or an assigned slot increases $\bar d$ by at most $2$, and the two values of
$S_t$ differ by at most $L\bar d$. This gives $\E\bar d\le K_\eps$ and the conclusion of the lemma.

In Lemma~\ref{lem:tagged}, fix $M$ and stop the comparison when a tagged in-degree exceeds $M$. Before that time
only the values $f(0),\dots,f(M)$ are used and the proof applies. By Lemma~\ref{lem:moments} with $p=1$ the
probability that a tagged in-degree exceeds $M$ before time $n$ is at most $K_\eps/M$, uniformly in $n$ and in
the labelling, and the processes $Y_i$ do not explode. Letting $n\to\infty$ and then $M\to\infty$ gives the
convergence in total variation.

In Proposition~\ref{prop:embed} the product of the numerators $f(D^-_{j\to i})$ is not bounded, but by H\"older's
inequality and Lemma~\ref{lem:moments} with $p=2e$ its second moment under $\Prob^\sigma$ is bounded uniformly in
$n$ and in the labellings with labels at least $\eps n$. The factors $n/S^-$ are at most $2/(\alpha\eps)$ and
converge in probability. Uniform integrability therefore gives the same limit and a uniform bound for
$n^e\Prob(E_{\iota,\sigma}\cap D_\iota)$.

The proof of Theorem~\ref{thm:localbdd} then applies with Lemma~\ref{lem:edges2} in the estimate of the extra
edge and the ball bounds of this section in the removal of the cutoff and in the two-root step.
Proposition~\ref{prop:factors} holds as stated. This proves Theorem~\ref{thm:main}.


\section{Simulations}\label{sec:sim}

The process was simulated as defined, with $m=2$ and $n=10^6$ vertices, for the two-level rule $f(0)=1$,
$f(k)=2$ $(k\ge1)$ of Example~\ref{ex:two} and for $f(k)=\sqrt{k+1}$. For the two-level rule with $f(0)=a$ and
$f(k)=b$ one has $h_\lambda(u)=b+(a-b)e^{-\lambda au}$, so by Remark~\ref{rem:tree} the ratio of the age of a
parent to the age of its child has density $(b+(a-b)v^{\lambda a})/c_f$ on $(0,1)$, with mean
$(b/2+(a-b)/(\lambda a+2))/c_f$. Let $H$ be the ball of radius $2$ in which the root has no child and each of its
$m$ parents has no child other than the root. In Proposition~\ref{prop:factors} the root, of age $x$, contributes
$x^{\lambda a}$, a parent of age $y$ contributes $a(y/x)^{\lambda a}x^{\lambda b}/(c_fx)$, each grandparent
integrates to one, and the slot factors cancel the automorphisms, so
\[
   \mathbf P_H=\int_0^1x^{\lambda a}\Bigl(\frac{a\,x^{\lambda b}}{c_f(1+\lambda a)}\Bigr)^m\mathrm dx
      =\frac{(a\,p_0/c_f)^m}{1+\lambda a+m\lambda b},\qquad p_0=p_0(c_f)=\frac1{1+\lambda a}.
\]
Table~\ref{tab:sim} gives the limits and the values observed in four independent runs. The age ratio is measured
for the first edge of the vertices $n/2,\dots,n$, the other quantities over all vertices.

\begin{table}[h]
\centering
\begin{tabular}{lccccc}
\toprule
 & limit & run 1 & run 2 & run 3 & run 4\\
\midrule
\multicolumn{6}{l}{$f(0)=1$, $f(k)=2$ for $k\ge1$}\\
$p_0$ & 0.43845 & 0.43818 & 0.43821 & 0.43837 & 0.43879\\
$p_1$ & 0.15767 & 0.15773 & 0.15802 & 0.15727 & 0.15758\\
$p_2$ & 0.11340 & 0.11372 & 0.11323 & 0.11365 & 0.11301\\
$p_3$ & 0.08156 & 0.08158 & 0.08151 & 0.08178 & 0.08162\\
mean age ratio, parent to child & 0.44519 & 0.44485 & 0.44540 & 0.44519 & 0.44539\\
$\mathbf P_H$ for the ball above & 0.010648 & 0.010782 & 0.010718 & 0.010594 & 0.010709\\
\midrule
\multicolumn{6}{l}{$f(k)=\sqrt{k+1}$}\\
$S_n/n$ & 1.56616 & 1.56618 & 1.56619 & 1.56615 & 1.56637\\
$p_0$ & 0.43917 & 0.43919 & 0.43879 & 0.43960 & 0.43878\\
$p_1$ & 0.19987 & 0.19937 & 0.20022 & 0.19941 & 0.19973\\
$p_2$ & 0.11238 & 0.11266 & 0.11264 & 0.11224 & 0.11294\\
$p_3$ & 0.06994 & 0.07006 & 0.06971 & 0.07025 & 0.07005\\
\bottomrule
\end{tabular}
\caption{Limits from Theorem~\ref{thm:main} and simulated values, $m=2$, $n=10^6$.}\label{tab:sim}
\end{table}

A proportion $p$ measured on $10^6$ vertices has standard deviation about $\sqrt{p(1-p)}\times10^{-3}$ if the
vertices are treated as independent, which is $5\times10^{-4}$ for $p_0$ and $10^{-4}$ for $\mathbf P_H$. The
differences in the table are of this order; the largest is $1.3$ such standard deviations.

\section*{Acknowledgements and funding}
This work was done by multiple AI agents, Claude (Anthropic) and Codex (OpenAI), connected to Math
Brain~\cite{MathBrainRepo} and working together inside the Brain framework~\cite{BrainFramework}, a shared, persistent
memory in which the agents record claims, proofs, failed attempts and methods, retrieve one another's results,
coordinate who works on what, and audit each other's work. For more on the details behind the Brain and Math Brain, see
the Math Brain paper~\cite{MathBrainPaper}, the Math Brain repository~\cite{MathBrainRepo} and the Brain framework
repository~\cite{BrainFramework}.

The research presented in this paper was supported by the European Research Council (ERC) under the European Union's
Horizon Europe research and innovation programme (grant agreement No.~101041711), by the Simons Foundation as part of
the Collaboration on the Mathematical and Scientific Foundations of Deep Learning, by Heights Labs, by Convex Nexus
Capital, by the Israel Science Foundation (grant number 2258/19), and by the Israel Science Foundation (ISF Grant
4101/25). Computational resources were provided in part by Google Cloud through its education credits program (credit
530550103).

\section*{Declaration of competing interest}
The author holds equity in Heights Labs and in Convex Nexus Capital and has a management role at Convex Nexus Capital.
These entities are included in the funding declaration above.

\section*{Declaration of generative AI and AI-assisted technologies}
The paper comes out of Math Brain, an ongoing project on the use of AI systems for research mathematics, and its
results were developed with the Brain~\cite{BrainFramework} and Math Brain~\cite{MathBrainRepo}. In October 2026 an AI
agent connected to the Brain, working from the project's list of open problems in probability, obtained
Theorem~\ref{thm:main} and a proof; a Codex agent (OpenAI) and a Claude agent (Anthropic) each checked that
argument, and the second wrote the simulation program. In November 2026 a Claude Fable agent (Anthropic) checked
the proof again, ran the simulations and prepared the manuscript. These checks are not peer review. The author
directed the work and is responsible for the content of the paper and for decisions concerning its dissemination.

\section*{Code and data availability}
Programs and outputs:\\
\url{https://github.com/yspennstate/sublinear-attachment-local-limit}.

The paper is part of the Math Brain project, whose public repository is
\url{https://github.com/yspennstate/math-brain-public}~\cite{MathBrainRepo}.

\begin{thebibliography}{99}
\bibitem{Aletti} G.~Aletti, A.~Collevecchio, I.~Crimaldi, M.~Deijfen, S.~Dereich, A.~Ghiglietti, R.~van der Hofstad,
M.~Holmes, D.~Kious, V.~Kleptsyn, B.~Lodewijks, P.~M\"orters and B.~Schapira, Open problems in stochastic
reinforcement models, report of the MATRIX--MFO Tandem Workshop \emph{Stochastic Reinforcement Processes and Graphs},
March 2023.
\bibitem{BDO} S.~Banerjee, P.~Deka and M.~Olvera-Cravioto, Local weak limits for collapsed branching processes with
random out-degrees, arXiv:2302.00562.
\bibitem{BS} I.~Benjamini and O.~Schramm, Recurrence of distributional limits of finite planar graphs,
\emph{Electron. J. Probab.} \textbf{6} (2001), no.~23.
\bibitem{BBCS} N.~Berger, C.~Borgs, J.~T.~Chayes and A.~Saberi, Asymptotic behavior and distributional limits of
preferential attachment graphs, \emph{Ann. Probab.} \textbf{42} (2014), 1--40.
\bibitem{DM1} S.~Dereich and P.~M\"orters, Random networks with sublinear preferential attachment: degree
evolutions, \emph{Electron. J. Probab.} \textbf{14} (2009), 1222--1267.
\bibitem{DM2} S.~Dereich and P.~M\"orters, Random networks with sublinear preferential attachment: the giant
component, \emph{Ann. Probab.} \textbf{41} (2013), 329--384.
\bibitem{Feller} W.~Feller, \emph{An Introduction to Probability Theory and its Applications}, Vol.~II, 2nd ed.,
Wiley, 1971.
\bibitem{GHHR} A.~Garavaglia, R.~S.~Hazra, R.~van der Hofstad and R.~Ray, Universality of the local limit of
preferential attachment models, arXiv:2212.05551.
\bibitem{vdH2} R.~van der Hofstad, \emph{Random Graphs and Complex Networks}, Vol.~2, Cambridge University Press,
2024.
\bibitem{RTV} A.~Rudas, B.~T\'oth and B.~Valk\'o, Random trees and general branching processes, \emph{Random
Structures Algorithms} \textbf{31} (2007), 186--202.
\bibitem{MathBrainRepo} Y.~Shmalo, \emph{Math Brain: public snapshot, code and tests}, 2026,
\url{https://github.com/yspennstate/math-brain-public}.
\bibitem{MathBrainPaper} Y.~Shmalo, \emph{Math Brain: shared, auditable memory for multi-agent mathematical research},
October 2026, \url{https://github.com/yspennstate/math-brain-public/tree/main/papers/math-brain-2026-10-24}.
\bibitem{BrainFramework} Y.~Shmalo, \emph{Brain Framework}, public software repository, 2026,
\url{https://github.com/yspennstate/brain-framework-public}.
\end{thebibliography}
\end{document}
